%% kapa: exam=1 round=28 year=2017 period=2 form=A subject=accounting subject_ko=회계학 seq=30 number=70
%% source: Q-Net 감정평가사 기출문제 (cst00309) artlSeq=5199603 file=extracted/1008/5199603/2017년 제28회 감정평가사 1차시험 문제지/2교시 A형.pdf
%% conversion: kordoc 4.24.51 -> tools/gen1.py (auto); review=proofread-sonnet; flags=none
\begin{question}{70}
충당부채와 우발부채에 관한 설명으로 옳지 않은 것은?
\begin{choices}
  \item 제3자와 연대하여 의무를 지는 경우에는 이행할 전체의무 중 제3자가 이행할 것으로 예상되는 부분을 우발부채로 인식한다.
  \item 충당부채로 인식되기 위해서는 과거사건의 결과로 현재의무가 존재하여야 한다.
  \item 충당부채를 현재가치로 평가할 때 할인율은 부채의 특유한 위험과 화폐의 시간 가치에 대한 현행 시장의 평가를 반영한 세전 이율을 적용한다.
  \item 충당부채와 관련하여 포괄손익계산서에 인식한 비용은 제3자의 변제와 관련하여 인식한 금액과 상계하여 표시할 수 있다.
  \item 과거에 우발부채로 처리하였다면 이후 충당부채의 인식조건을 충족하더라도 재무제표의 신뢰성 제고를 위해서 충당부채로 인식하지 않는다.
\end{choices}
\end{question}
\answer{5}
